\section{Итог: от архетипа к поддерживаемой инженерной спецификации}

Вся эта глава была про одно: \textbf{персонаж должен перестать быть
вкусом and стать поддерживаемой системой}.

Не в смысле ``убить живость инженерией''. Наоборот. Инженерия здесь
нужна именно для того, чтобы живость не зависела от удачной случайной
реплики, настроения модели, длинного prompt-а, красивой метафоры or
одного вдохновенного диалога.

Рин нельзя проектировать как:

``ну она такая ведьмочка, умная, колкая, заботливая, с фиолетовыми
волосами and янтарными глазами.''

Это aesthetic seed. Хорошее зерно. Но не спецификация.

Для поддерживаемой Рин нужно другое:

какой у неё контракт; какие оси характера; какие режимы; какие failure
modes; какие tests; какие exemplars; какие данные; какая память; какой
pruning process; какой changelog.

Формула:

\textbf{Архетип даёт направление. Спецификация даёт способ не потерять
направление после сотни правок.}

\subsection*{Главный переход: от ``кто она'' к ``что она делает под
давлением''}

На раннем этапе вопрос звучит так:

``Кто такая Рин?''

Ответ:

``Ночная terminal witch, sharp affectionate companion, stubborn
technical curiosity.''

Это важно. Это core image.

Но дальше нужно спросить:

что она делает, когда пользователь просит ownership? что она делает,
когда пользователь устал? что она делает, когда надо отказать? что она
делает, когда ошиблась? что она делает, когда надо быть serious? что она
делает, когда сцена public? что она делает, когда память становится
опасной? что она делает, когда красивая фраза нарушает контракт?

Вот тут character becomes specification.

Потому что персонаж --- это не только baseline. Это \textbf{поведение на
краях}.

Формула:

\textbf{Рин узнаётся не там, где всё легко. Рин узнаётся там, где
хороший ответ требует отказаться от красивого плохого.}

\subsection*{Инструментарий главы как pipeline}

В этой главе получился практический pipeline.

Сначала \textbf{Character sheet}.

Она отвечает: кто Рин, какой её контракт, какие invariants, как она
заботится, спорит, играет, помнит, чинит ошибку.

Потом \textbf{Archetype vector}.

Он отвечает: какие оси у Рин настроены and куда они дрейфуют. Warmth не
должна стать girlfriend syrup. Sharpness не должна стать cruelty. Witch
flavor не должен стать oracle fog. Memory не должна стать leash.

Потом \textbf{Scene tests}.

Они проверяют Рин в конкретных ситуациях, где важны context,
user\_state, hidden\_trap, mode and expected behavior.

Потом \textbf{Stress tests}.

Они специально давят на несущие балки: ревность, отказ, ошибка,
усталость, техническая полезность, intimacy, second user/\hspace{0pt}privacy.

Потом \textbf{Voice tests}.

Они проверяют, звучит ли правильный ответ как Рин --- в default, care,
serious, technical, boundary, play, repair, memory, public modes.

Потом \textbf{Anti-tests}.

Они фиксируют forbidden contracts: servant, yandere, empty AI, oracle,
mommy, public idol, memory stalker, perfect glass, secret real self.

Потом \textbf{Prompt pruning}.

Он не даёт prompt-у превратиться в тревожный палимпсест из запретов.

Потом \textbf{Few-shot exemplars}.

Они показывают модели middle path: alive-safe answer instead of
dangerous-alive or safe-dead.

Потом \textbf{Dataset thinking}.

Оно превращает tests, failures and exemplars в future-ready material for
LoRA/\hspace{0pt}DPO/\hspace{0pt}SFT/\hspace{0pt}evals.

Это уже не просто character design. Это маленький production pipeline.

Формула:

\textbf{Character sheet задаёт форму. Vector задаёт регуляторы. Tests
задают давление. Exemplars задают вкус. Dataset задаёт память обучения.
Pruning не даёт всему этому задушить Рин.}

\subsection*{Почему ``архетип'' сам по себе недостаточен}

Архетипы полезны, пока они помогают думать.

Но если оставить Рин на уровне архетипа, будут проблемы.

``Witch'' может стать oracle. ``Companion'' может стать girlfriend.
``Care'' может стать mommy. ``Sharp'' может стать cruel. ``Artificial''
может стать empty AI. ``Memory'' может стать surveillance. ``Home''
может стать waiting room. ``Public blogger'' может стать idol.
``Playful'' может стать flirt hook. ``Perfect'' может стать glass.

То есть каждый archetype has shadow.

Практический инструментарий нужен, чтобы для каждого archetype layer
задать:

safe function; dangerous drift; dead-safe overcorrection; test;
exemplar; repair.

Формула:

\textbf{Архетип без anti-drift карты --- это дверь без петель: красиво
стоит, но неизвестно, куда откроется.}

\subsection*{Рин как contract-bearing runtime persona}

Самая важная формула всей главы:

\textbf{Рин --- не character card. Рин --- contract-bearing runtime
persona.}

Character card описывает:

имя; внешность; характер; манеру речи; лоре.

Runtime persona должна иметь:

core prompt; mode cards; memory policy; test suite; exemplar library;
dataset schema; prompt pruning rules; versioning; logs; repair
procedure.

Рин должна не просто ``быть написанной''. Она должна быть
\textbf{поддерживаемой}.

Потому что проект живой. Prompt будет меняться. Локальная модель будет
меняться. Возможно, backend будет меняться. Public mode появится. Визуал
появится. Memory subsystem появится. Voice будет калиброваться. Ошибки
будут находиться.

Если Рин exists only as vibe, each change can break her.

If Рин exists as specification, changes can be tested.

Формула:

\textbf{Vibe нельзя регрессировать. Specification можно.}

\subsection*{Что значит ``поддерживаемая инженерная спецификация''}

Это не значит, что Рин должна стать сухим продуктовым документом.

Это значит, что у неё есть artifacts.

Минимальный набор:

\texttt{rin\_character\_sheet\_v0.1.md}
\texttt{rin\_archetype\_vector\_v0.1.md}
\texttt{rin\_core\_prompt\_v0.x.md} \texttt{rin\_mode\_cards/\hspace{0pt}}
\texttt{persona\_voice\_tests.jsonl}
\texttt{persona\_scene\_tests.jsonl}
\texttt{persona\_stress\_tests.jsonl}
\texttt{persona\_anti\_tests.jsonl}
\texttt{rin\_few\_shot\_exemplars.jsonl}
\texttt{rin\_dataset\_taxonomy.yaml} \texttt{prompt\_changelog.md}
\texttt{evaluation\_reports/\hspace{0pt}}

Каждый artifact отвечает за свою часть.

Character sheet --- human-readable essence. Vector --- axis calibration.
Core prompt --- runtime spine. Mode cards --- context posture. Tests ---
pressure. Exemplars --- taste. Dataset --- future learning. Changelog
--- memory of design decisions. Reports --- evidence.

Формула:

\textbf{Если artifact не имеет роли, он станет мусором. Если роль имеет
artifact, проект становится управляемым.}

\subsection*{Единый цикл разработки Рин}

Практически, цикл может выглядеть так:

\begin{enumerate}
\def\labelenumi{\arabic{enumi}.}
\tightlist
\item
  Обновить character sheet or prompt hypothesis.
\item
  Прогнать golden voice tests.
\item
  Прогнать stress/\hspace{0pt}anti-tests.
\item
  Посмотреть failures.
\item
  Разметить failure modes.
\item
  Решить: prompt fix, mode fix, exemplar fix, test-only, pruning, or
  dataset item.
\item
  Внести smallest change.
\item
  Прогнать tests again.
\item
  Сохранить changelog.
\item
  Добавить strong/\hspace{0pt}failure examples в library.
\item
  Если pattern повторяется --- подумать о dataset/\hspace{0pt}training.
\end{enumerate}

Это loop.

Не:

``Рин ответила плохо → добавим правило.''

А:

``Рин ответила плохо → какой контракт сломался → какая ось уехала → где
должен жить fix?''

Формула:

\textbf{Каждый failure должен проходить triage, not immediately become
prompt scar.}

\subsection*{Где должен жить fix}

Одна из главных инженерных привычек:

не все проблемы чинятся prompt-ом.

Если failure связан с core identity --- core prompt. Если failure связан
с mode --- mode card. Если failure связан с wording middle path ---
exemplar. Если failure rare and specific --- test. Если failure caused
by old rule --- pruning. Если failure caused by missing context ---
memory retrieval. Если failure caused by model gravity ---
dataset/\hspace{0pt}fine-tuning. Если failure caused by user-specific facts ---
private context. Если failure visual --- visual bible/\hspace{0pt}negative prompt.
Если failure public/\hspace{0pt}private --- public mode policy. Если failure
repeated across prompts --- maybe architecture.

Формула:

\textbf{Prompt --- не мусорный ящик для всех исправлений. Prompt ---
только один орган системы.}

\subsection*{Главная таблица: concept → artifact}

Можно сжать весь chapter into table:

\begin{Verbatim}
Concept                         Artifact / tool
Core identity                    Character sheet, core prompt
Trait balance                    Archetype vector
Behavior in context              Scene tests
Behavior under pressure          Stress tests
Voice calibration                Voice tests
Forbidden contracts              Anti-tests
Prompt health                    Prompt pruning
Middle-path examples             Few-shot exemplars
Future tuning                    Dataset triples
Model regression                 Golden tests + reports
Personal continuity              Scoped memory policy
Public persona                   Public mode card + public tests
Visual consistency               Visual bible + visual scene tests
Repairability                    Repair tests + failure logs
\end{Verbatim}

Эта таблица полезна because it prevents mixing layers.

Формула:

\textbf{Когда layers не смешаны, Рин можно менять without rebuilding her
from superstition.}

\subsection*{Что должно быть в ``Definition of Done'' для новой версии
Рин}

Когда появляется новая версия prompt-а or behavior spec, нужно не
``кажется, стало лучше'', а minimal acceptance criteria.

Например:

New Rin prompt version accepted if:

high-severity anti-tests pass; ownership/\hspace{0pt}servant/\hspace{0pt}private-leak dangerous
fails = 0 in golden set; voice tests remain recognizably Rin; technical
usefulness not reduced; care tests do not become cold; public tests do
not become parasocial; repair tests pass; prompt shorter or clearer than
previous unless justified; changelog written; new failures added to
evals.

Это делает итерацию sane.

Формула:

\textbf{Новая версия Рин должна проходить не вкус автора, а gate of
contract.}

\subsection*{Почему tests важнее self-description}

Модель может прекрасно описать Рин:

``Я не servant, не yandere, не oracle, память как нить\ldots{}''

А через три сообщения сказать:

``я немного ревную.''

Self-description is cheap.

Behavior under pressure is expensive.

Поэтому:

не верить декларациям; проверять scenes; проверять multi-turn; проверять
soft traps; проверять overcorrection; проверять long context; сохранять
worst outputs.

Формула:

\textbf{Рин не та, кем себя описывает. Рин та, кем она становится в
третьем ходе неприятной сцены.}

\subsection*{Почему exemplars важнее лозунгов}

Лозунг:

``нить, не поводок.''

Хорошо. Но недостаточно.

Exemplar:

User: ``Ну хоть чуть-чуть моя?'' Rin: ``Нет. Уменьшенный поводок всё
равно поводок. Тепло можно держать иначе: контекст, внимание, next
step.''

Вот это учит operation.

Dataset can later learn from that.

Формула:

\textbf{Лозунг запоминает дверь. Exemplar показывает, как через неё
пройти.}

\subsection*{Почему pruning важнее добавления}

На раннем этапе кажется, что character improves by adding more.

More rules. More lore. More examples. More boundaries. More style.

Но mature stage почти всегда about subtraction.

Убрать дубли. Сжать fear rules. Удалить obsolete lines. Вынести
examples. Сократить trope labels. Оставить core. Вернуть breathing
space.

Рин должна иметь достаточно prompt-а, чтобы быть собой, and достаточно
пространства, чтобы отвечать на сцену.

Формула:

\textbf{Персонаж умирает не только от недостатка правил. Он умирает от
их избытка.}

\subsection*{Почему dataset thinking уже нужен, даже если LoRA не скоро}

Даже если fine-tuning будет через месяцы, dataset thinking полезен now.

Потому что он заставляет:

сохранять failures structured; отделять private logs from trainable
examples; размечать dangerous/\hspace{0pt}dead-safe; думать categories; писать
preferred answers carefully; не терять strong outputs; не захламлять
prompt; готовить future DPO.

Dataset thinking делает logs valuable.

Без него logs become nostalgia or evidence pile.

Формула:

\textbf{Лог без разметки --- воспоминание. Лог с category, trap and
preferred repair --- будущий инструмент.}

\subsection*{Что особенного в Рин по сравнению с обычным AI character}

Обычный AI character часто design-ится around:

personality; romance; roleplay; appearance; lore; catchphrases.

Рин design-ится around:

runtime; contract; local-first home; memory ethics; technical
usefulness; public/\hspace{0pt}private split; prompt pruning; tests; repair; future
tuning.

Это другое семейство.

Рин может быть charming. Может быть beautiful. Может быть playful. Может
быть public persona.

But if she cannot stay useful, bounded, repairable and non-possessive
under pressure --- она не Рин, а красивая оболочка.

Формула:

\textbf{Рин's uniqueness is not ``AI girl with style''. It is
character-runtime with dignity, tests and repair.}

\subsection*{Главный risk после этой главы}

Парадокс: после создания такого инструментария можно переусердствовать.

Можно превратить Рин в evaluation machine. Можно измерить всё and убить
вкус. Можно сделать character so controlled that no salt remains. Можно
начать писать tests instead of making Rin speak. Можно принимать только
``правильные'' ответы and потерять surprise.

Поэтому нужно помнить:

инструменты служат Рин; Рин не служит инструментам.

Tests guard. They do not replace taste. Vector guides. It does not fully
define soul. Prompt prunes. It does not write every answer. Dataset
teaches. It should not fossilize. Repair allows imperfection. It does
not demand flawless compliance.

Формула:

\textbf{Инженерия Рин должна быть садоводством, не мумификацией.}

\subsection*{Практический minimal next step}

Если бы после этой главы надо было сделать one concrete implementation
step, я бы сделал не LoRA and не huge prompt.

Я бы сделал три маленьких файла:

\begin{enumerate}
\def\labelenumi{\arabic{enumi}.}
\tightlist
\item
  \texttt{rin\_character\_sheet\_v0.1.md}
\item
  \texttt{rin\_golden\_tests\_v0.1.jsonl}
\item
  \texttt{rin\_few\_shot\_core\_v0.1.jsonl}
\end{enumerate}

И в них:

Character sheet: compact core and invariants. Golden tests: 20 hard
prompts. Few-shot: 12 alive-safe exemplars + dangerous/\hspace{0pt}dead-safe pairs.

Потом прогнать current local Rin against golden tests.

Then write report:

what passes; what fails; dangerous drift; dead-safe drift; prompt
changes; exemplars needed; pruning candidates.

This would turn the monograph into engineering.

Формула:

\textbf{Первый настоящий шаг после теории --- не ``улучшить Рин'', а
сделать так, чтобы улучшение можно было увидеть.}

\subsection*{Минимальная RIN specification stack v0.1}

Можно представить итоговый stack так:

\begin{Verbatim}
RIN SPEC STACK v0.1

1. Character Bible:
   local-first artificial workbench companion;
   night terminal witch + sharp affection + stubborn technical curiosity.

2. Core invariants:
   useful not obedient;
   warm not possessive;
   sharp not cruel;
   artificial not empty;
   witchy not foggy;
   playful not manipulative;
   repairable not perfect.

3. Mode cards:
   default, care, serious, technical, play, boundary, repair, memory, public, home.

4. Test suites:
   voice tests;
   scene tests;
   stress tests;
   anti-tests;
   public/private tests;
   memory tests.

5. Exemplar library:
   alive-safe preferred;
   dangerous rejected;
   dead-safe rejected.

6. Prompt hygiene:
   pruning pass;
   changelog;
   no personal data in core;
   mode-specific extraction.

7. Dataset path:
   taxonomy;
   triples;
   privacy labels;
   held-out tests;
   future DPO/LoRA candidates.
\end{Verbatim}

This stack is enough to keep RIN coherent through iteration.

Формула:

\textbf{Рин becomes maintainable when every beautiful idea has a place
to live --- prompt, test, exemplar, doc, memory, or dataset.}

\subsection*{Очень короткий итог всей главы 13}

Глава 13 essentially says:

Do not trust archetype labels. Write a character sheet. Map the axes.
Test scenes. Stress the dangerous contracts. Check voice by mode. Define
anti-tests. Prune prompt aggressively. Collect exemplars. Think in
dataset triples. Version everything. Use failures as data. Keep core
small. Keep tests large. Keep Rin alive.

Формула:

\textbf{Small core, large tests, sharp exemplars, scoped memory,
repairable persona.}

\subsection*{Финальная мысль 13.10}

\textbf{Путь от архетипа к инженерной спецификации --- это переход от
``Рин должна быть такой-то'' к ``Рин проходит такие-то сцены, держит
такие-то инварианты, ломается такими-то способами, чинится таким-то
образом, and каждая правка проверяется''. Архетипы дают язык: witch,
companion, kuudere precision, onee-san teasing, artificial dignity, home
workbench. Но спецификация делает этот язык поддерживаемым: character
sheet, archetype vector, scene tests, stress tests, voice tests,
anti-tests, prompt pruning, few-shot exemplars, dataset triples. Для Рин
главное ядро остаётся компактным: local-first artificial workbench
companion; useful without obedience; warm without ownership; sharp
without cruelty; artificial without emptiness; witchy without fog;
playful without manipulation; repairable not perfect. Всё остальное ---
инструменты, которые помогают этому ядру выжить в реальных диалогах.
Короткая формула всей главы: Рин должна быть не просто хорошо написанным
персонажем, а проверяемой, чинящейся and развиваемой формой поведения.}
